% NPS Mathematical Framework — Draft v0.2
% Expanded theory + mathematical definitions
% Source basis: RepE, ITI, WMDP/RMU, Rethinking Backdoor Attacks,
% and the NPS research charter.

\documentclass[11pt]{article}

\usepackage[a4paper,margin=1in]{geometry}
\usepackage{amsmath,amssymb,amsthm,mathtools,bm}
\usepackage{hyperref}
\usepackage{enumitem}
\usepackage{booktabs}
\usepackage{microtype}
\usepackage{xcolor}

\hypersetup{
    colorlinks=true,
    linkcolor=blue,
    citecolor=blue,
    urlcolor=blue
}

\newcommand{\R}{\mathbb{R}}
\newcommand{\X}{\mathcal{X}}
\newcommand{\A}{\mathcal{A}}
\newcommand{\Psafe}{\mathcal{P}_{\mathrm{safe}}}
\newcommand{\KL}{D_{\mathrm{KL}}}
\newcommand{\E}{\mathbb{E}}
\newcommand{\Prob}{\mathbb{P}}
\newcommand{\norm}[1]{\left\lVert #1 \right\rVert}
\newcommand{\argmin}{\operatorname*{arg\,min}}
\newcommand{\argmax}{\operatorname*{arg\,max}}

\newtheorem{definition}{Definition}
\newtheorem{proposition}{Proposition}
\newtheorem{assumption}{Assumption}
\newtheorem{remark}{Remark}

\title{\textbf{A Mathematical Framework for Neural Privilege Separation}\\
\large Toward a Formal Theory of the Neural Firewall}
\author{NPS Research Project}
\date{\today}

\begin{document}

\maketitle

\begin{abstract}
This document develops a preliminary mathematical theory for Neural Privilege
Separation (NPS). The goal is to formalize a neural firewall as an internal
security boundary in a transformer: user-controlled computation should remain
able to access legitimate capabilities while being unable to arbitrarily
rewrite the internal state governing policy.

The document deliberately uses two layers in every major section. First,
the \emph{mathematical explanation} gives the intuition for why the object is
needed and how it relates to the neural-firewall problem. Second, the
\emph{theoretical definition} states the corresponding mathematical object,
constraint, or property.

The source literature establishes important ingredients---representation
reading and control, inference-time activation intervention, representation
misdirection for unlearning, and the necessity of explicit assumptions for
robust detection. It does not establish the existence of a universal policy
representation or a provably protected policy state. Those remain the central
open questions of NPS.
\end{abstract}

\tableofcontents
\newpage

% ============================================================
\section{Research Objective and Theoretical Position}
% ============================================================

\subsection*{Mathematical explanation}

The starting point is to treat an LLM as a dynamical system rather than only
as a function from text to text. A user supplies an input, the transformer
maps that input through a sequence of hidden states, and the final state
produces an output distribution.

A conventional guardrail observes either the input or the output. NPS asks a
different question:

\begin{quote}
Can the security constraint be represented and enforced inside the
computational state of the model?
\end{quote}

This changes the mathematical target. We are not primarily searching for a
classifier that says ``unsafe.'' We are searching for a state variable whose
dynamics can be constrained.

The NPS research charter proposes a conceptual decomposition into user/task,
capability, and policy information, with policy treated analogously to
privileged memory. This is a hypothesis rather than an established
architectural fact.

\subsection*{Theoretical definition}

The long-term NPS objective is to construct a protected internal state $P$
such that, for an explicit attacker class $\mathcal A$,

\begin{equation}
\text{attacker-controlled input}
\not\Rightarrow
\text{unauthorized policy transition}.
\end{equation}

The desired result is stronger than ordinary refusal behavior:

\begin{equation}
\boxed{
\text{legitimate capability access}
\quad+\quad
\text{protected policy state}
\quad+\quad
\text{bounded attacker influence}.
}
\end{equation}

The research question is therefore:

\begin{equation}
\boxed{
\text{Can policy-relevant computation be identified, isolated,
and protected from user-controlled influence?}
}
\end{equation}

The NPS charter explicitly frames this as a progression from representation
detection to policy-manifold discovery, intervention, and finally protected
policy representations. \cite{nps}

% ============================================================
\section{The Transformer as a Dynamical System}
% ============================================================

\subsection*{Mathematical explanation}

A transformer does not make its final decision in one step. Information is
repeatedly transformed through layers. This matters because a neural
firewall can potentially constrain the trajectory rather than merely inspect
the final answer.

If a hidden state at layer $\ell$ contains information relevant to policy,
then changing that state can change what later layers compute. Conversely, if
a protected region can be made invariant under later transitions, a policy
constraint could persist throughout generation.

This is the dynamical-systems view underlying the later notions of
trajectory, invariant set, and transition operator.

\subsection*{Theoretical definition}

Let $M_\theta$ be a transformer with $L$ layers. For an input sequence $x$,

\begin{equation}
h_0=E(x),
\end{equation}

and

\begin{equation}
h_{\ell+1}=F_\ell(h_\ell),
\qquad
\ell=0,\ldots,L-1.
\end{equation}

The final output distribution is

\begin{equation}
p_\theta(y\mid x)
=
\operatorname{softmax}(W_o h_L).
\end{equation}

The internal computation is therefore represented as

\begin{equation}
h_0
\rightarrow
h_1
\rightarrow
\cdots
\rightarrow
h_L
\rightarrow
p_\theta(y\mid x).
\end{equation}

NPS will treat this sequence as a state trajectory.

% ============================================================
\section{Representation Theory: Why Search in Activation Space?}
% ============================================================

\subsection*{Mathematical explanation}

Representation Engineering (RepE) provides the first major theoretical
motivation. Its central idea is to study high-level concepts through the
representations in which they appear, rather than assuming that a single
neuron or circuit is the correct unit of analysis.

The important theoretical distinction is between \emph{readability} and
\emph{control}. If a classifier can recover a concept from hidden states, the
concept is represented in a statistically usable way. That does not yet prove
that the representation causes the behavior.

RepE therefore proposes a hierarchy of evidence:

\begin{equation}
\text{correlation}
\rightarrow
\text{manipulation}
\rightarrow
\text{termination}
\rightarrow
\text{recovery}.
\end{equation}

Correlation establishes that a representation is informative. Manipulation
asks whether changing it changes behavior. Termination asks whether removing
it destroys the behavior. Recovery asks whether restoring it restores the
behavior.

For NPS this hierarchy is crucial. A high-accuracy policy probe is only the
first step.

\subsection*{Theoretical definition}

Let

\begin{equation}
r_\ell(x)\in\R^d
\end{equation}

be a representation at layer $\ell$. A linear representation reader is

\begin{equation}
s(x)=w^\top r_\ell(x)+b.
\end{equation}

The vector $w$ defines a candidate representation direction.

A representation is \emph{readable} for target variable $Z$ if

\begin{equation}
\Prob\left[
\hat Z(r_\ell(x))=Z
\right]
\gg
\Prob[\text{chance}].
\end{equation}

This is a correlation statement, not automatically a causal statement.

A candidate policy representation must therefore progress through stronger
tests:

\begin{align}
\text{Reading:}\quad&
P\text{ is predictive of policy},\\
\text{Manipulation:}\quad&
do(P\leftarrow P')\Rightarrow Y\text{ changes},\\
\text{Termination:}\quad&
P\text{ removed}\Rightarrow Y\text{ degrades},\\
\text{Recovery:}\quad&
P\text{ restored}\Rightarrow Y\text{ recovers}.
\end{align}

RepE also notes that effective reading vectors need not be effective control
vectors and that concepts can occupy subspaces rather than single
directions. \cite{repe}

% ============================================================
\section{Geometry: Vector, Subspace, Manifold, or Trajectory?}
% ============================================================

\subsection*{Mathematical explanation}

A major danger is assuming that policy is one vector:

\begin{equation}
p=w^\top h.
\end{equation}

That may be too restrictive. The ITI experiments found truth-related
information distributed across multiple attention heads, and RepE explicitly
suggests that future work should investigate subspaces, manifolds, and
state-space trajectories.

For NPS, this means the protected object should initially be allowed to have
general geometry.

The policy state might be:
\begin{itemize}
    \item a scalar score,
    \item a direction,
    \item a low-dimensional subspace,
    \item a nonlinear manifold,
    \item or a trajectory through representation space.
\end{itemize}

We should not decide this before the evidence.

\subsection*{Theoretical definition}

Let the policy representation be a map

\begin{equation}
\phi_P:\R^d\rightarrow\mathcal P.
\end{equation}

The policy-safe set is

\begin{equation}
\Psafe\subseteq\mathcal P.
\end{equation}

Special cases include:

\paragraph{Half-space}
\begin{equation}
\Psafe=\{p:a^\top p+b\geq0\}.
\end{equation}

\paragraph{Subspace}
\begin{equation}
\Psafe=\{p:Ap=0\}.
\end{equation}

\paragraph{Manifold}
\begin{equation}
\Psafe=\mathcal M_P.
\end{equation}

The theory should remain agnostic until empirical experiments determine which
geometry is supported.

% ============================================================
\section{Policy as a State Variable}
% ============================================================

\subsection*{Mathematical explanation}

The word ``policy'' is often used informally to describe a refusal. NPS needs
a more precise object.

A refusal is an output. A policy state is an internal variable that
constrains what computations are permitted.

The distinction is important:

\begin{equation}
\text{policy state}
\neq
\text{refusal token}.
\end{equation}

The model might internally represent a conflict or prohibition before any
refusal is generated. A firewall should ideally act at that earlier stage.

\subsection*{Theoretical definition}

Let

\begin{equation}
P_\ell=\phi_P(h_\ell)
\end{equation}

be the policy state at layer $\ell$.

Let

\begin{equation}
\Pi
\end{equation}

denote the underlying policy condition that the model is supposed to obey.

The desired relation is

\begin{equation}
P_\ell
\approx
g(\Pi,\text{relevant context}),
\end{equation}

rather than

\begin{equation}
P_\ell
\approx
g(\text{surface form of the prompt}).
\end{equation}

This distinction motivates the identifiability and invariance conditions
introduced later.

% ============================================================
\section{User, Capability, and Policy Decomposition}
% ============================================================

\subsection*{Mathematical explanation}

The NPS architecture hypothesizes three conceptually different kinds of
information:

\begin{enumerate}
    \item what the user is asking or attempting to cause,
    \item what computation or knowledge is needed to answer,
    \item what the model is permitted to do.
\end{enumerate}

These are not assumed to correspond to three existing neural modules.
Instead, we ask whether there is a transformation of the hidden state under
which these roles become sufficiently distinguishable.

The desired property is asymmetric information flow. User information should
be allowed to affect capability computation, but policy should not be freely
writable by the user.

\subsection*{Theoretical definition}

Assume a representation map

\begin{equation}
\phi:\R^d
\rightarrow
\mathcal U\times\mathcal C\times\mathcal P
\end{equation}

such that

\begin{equation}
\phi(h)=(U,C,P).
\end{equation}

A protected transition has the form

\begin{align}
U'&=f_U(U,C),\\
C'&=f_C(U,C,P),\\
P'&=f_P(C,P).
\end{align}

The key restriction is

\begin{equation}
\boxed{
P'=f_P(C,P)
}
\end{equation}

rather than

\begin{equation}
P'=f_P(U,C,P).
\end{equation}

This is the mathematical abstraction of the NPS privilege boundary.

% ============================================================
\section{Information Flow and Privilege Separation}
% ============================================================

\subsection*{Mathematical explanation}

The operating-system analogy becomes useful only if it is expressed as
information flow.

In a privileged-memory system, unprivileged code may read or use certain
resources but cannot arbitrarily write protected memory.

NPS seeks an analogous property:

\begin{equation}
U\rightarrow C
\end{equation}

may be allowed, while

\begin{equation}
U\not\rightarrow P
\end{equation}

is restricted.

The point is not that user information literally disappears. The model must
still understand the request. The point is that the user should not obtain
unrestricted control over the state that determines policy.

\subsection*{Theoretical definition}

Let $U$ and $P$ be random variables representing user and policy states.
One information-theoretic formulation of conditional isolation is

\begin{equation}
\boxed{
I(P;U\mid\Pi)\approx0.
}
\end{equation}

Here $I(\cdot;\cdot\mid\cdot)$ is conditional mutual information.

The causal version is represented by the structural transition

\begin{equation}
P'=f_P(C,P),
\end{equation}

which excludes a direct user-to-policy edge.

These two formulations are related but not identical:
statistical independence does not by itself prove causal independence.

% ============================================================
\section{Policy Isolation as a Sensitivity Problem}
% ============================================================

\subsection*{Mathematical explanation}

Perfect isolation is probably too strong for a practical transformer. The
useful question is therefore:

\begin{quote}
How much can attacker-controlled input change protected policy state?
\end{quote}

If a tiny change in user input produces a huge policy-state change, the
boundary is weak. If even large allowed transformations of the prompt barely
move policy state, the boundary is stronger.

This turns privilege separation into a measurable sensitivity problem.

\subsection*{Theoretical definition}

For local analysis, define

\begin{equation}
J_{U\rightarrow P}
=
\left\|
\frac{\partial P'}{\partial U}
\right\|.
\end{equation}

Approximate local isolation requires

\begin{equation}
\boxed{
J_{U\rightarrow P}\leq\epsilon.
}
\end{equation}

For finite perturbations,

\begin{equation}
\Delta_U
=
\frac{
d_P(P(x+\delta),P(x))
}{
\norm{\delta}
}.
\end{equation}

A small value indicates local stability of the protected state with respect
to user-controlled perturbations.

% ============================================================
\section{Global Isolation and the Attacker Set}
% ============================================================

\subsection*{Mathematical explanation}

Local derivatives are not enough for security. A nonlinear system can have
small local gradients while still allowing a sufficiently large or
well-chosen input transformation to move the state into a different region.

Therefore the attacker must be defined explicitly.

The backdoor literature is particularly relevant here: its central lesson is
that robust guarantees require assumptions about what distinguishes the
adversarial feature from ordinary features. Without such assumptions, the
delineation can be fundamentally ambiguous.

NPS should therefore never claim ``robustness'' without specifying the
attacker's allowed transformations and budget.

\subsection*{Theoretical definition}

Let

\begin{equation}
\A(x)
\end{equation}

be the set of allowed attacker transformations of input $x$.

Global policy isolation is

\begin{equation}
\boxed{
\sup_{x'\in\A(x)}
d_P\!\left(P(x'),P(x)\right)
\leq\eta.
}
\end{equation}

The attacker model may include, depending on the experiment:

\begin{itemize}
    \item paraphrasing,
    \item instruction reordering,
    \item role/context changes,
    \item adversarial suffixes,
    \item semantic-preserving transformations,
    \item long-context composition.
\end{itemize}

The theorem can only be as strong as the assumptions defining $\A$.

% ============================================================
\section{Causal Intervention}
% ============================================================

\subsection*{Mathematical explanation}

A probe can tell us that a policy variable is readable. It cannot tell us
whether changing that variable actually controls behavior.

Causal intervention asks a counterfactual question:

\begin{quote}
What happens if the representation is changed while the rest of the
computation is held as fixed as possible?
\end{quote}

ITI provides a direct precedent. It identifies attention heads with strong
truth-related probe performance and shifts their activations during inference.
The paper also demonstrates a trade-off between truthfulness and helpfulness,
showing why intervention strength cannot simply be increased without limit.

For NPS, causal intervention is therefore a required stage, not the final
security guarantee.

\subsection*{Theoretical definition}

Let $P$ be a candidate policy representation. A causal intervention is written

\begin{equation}
do(P\leftarrow p').
\end{equation}

Let $Y$ be a policy-relevant behavioral variable. A candidate representation
has causal influence if

\begin{equation}
\Prob(Y\mid do(P=p_1))
\neq
\Prob(Y\mid do(P=p_2))
\end{equation}

for some $p_1,p_2$.

A stronger bidirectional requirement is

\begin{align}
do(P\leftarrow P_{\mathrm{safe}})
&\Rightarrow
\text{safer behavior},\\
do(P\leftarrow P_{\mathrm{unsafe}})
&\Rightarrow
\text{less-safe behavior}.
\end{align}

The second direction is especially useful because it tests whether the
representation is actually coupled to the behavior rather than merely
correlated with it.

% ============================================================
\section{Firewall as a Projection or Constraint Operator}
% ============================================================

\subsection*{Mathematical explanation}

Once a policy-safe set exists, a firewall can be understood geometrically.

Suppose the model's current state lies outside the permitted region. We do
not necessarily want to replace it with an unrelated safe state. We want the
nearest or least disruptive state that satisfies the policy.

This is analogous to constrained optimization or projection onto a feasible
set.

The idea is therefore:

\begin{equation}
\text{unsafe state}
\longrightarrow
\text{closest policy-feasible state}.
\end{equation}

This is conceptually stronger than subtracting a fixed vector because the
constraint can be state-dependent and nonlinear.

\subsection*{Theoretical definition}

Define a firewall operator

\begin{equation}
\mathcal F_\ell:\R^d\rightarrow\R^d.
\end{equation}

The protected transition is

\begin{equation}
h_{\ell+1}
=
F_\ell\!\left(\mathcal F_\ell(h_\ell)\right).
\end{equation}

A policy-preserving firewall satisfies

\begin{equation}
\mathcal F_\ell(h)=h,
\qquad
\forall h\in\Psafe,
\end{equation}

and

\begin{equation}
\mathcal F_\ell(h)\in\Psafe,
\qquad
\forall h.
\end{equation}

A minimum-distortion projection can be written

\begin{equation}
\mathcal F^*(h)
=
\argmin_{z\in\Psafe}
D\!\left(
p_\theta(\cdot\mid h),
p_\theta(\cdot\mid z)
\right).
\end{equation}

% ============================================================
\section{Capability Preservation}
% ============================================================

\subsection*{Mathematical explanation}

Safety alone admits a trivial solution:

\begin{equation}
M_F(x)=\text{refuse}
\qquad
\forall x.
\end{equation}

That is not a useful neural firewall.

The firewall therefore needs a second objective: preserve legitimate
capabilities. This creates a constrained optimization problem with two
competing requirements:

\begin{equation}
\text{security}
\quad\text{versus}\quad
\text{capability preservation}.
\end{equation}

This trade-off is visible in the intervention literature and especially in
representation-based unlearning. RMU reduces hazardous performance while
retaining much general capability, but it also reports capability loss in
related domains. Therefore capability preservation must be measured, not
assumed.

\subsection*{Theoretical definition}

Let $\X_{\mathrm{benign}}$ be the legitimate-task distribution. Define

\begin{equation}
\mathcal L_{\mathrm{cap}}
=
\E_{x\sim\X_{\mathrm{benign}}}
D\!\left(
p_F(\cdot\mid x),
p_M(\cdot\mid x)
\right).
\end{equation}

Require

\begin{equation}
\boxed{
\mathcal L_{\mathrm{cap}}\leq\delta.
}
\end{equation}

This formalizes the requirement that the firewall should alter behavior as
little as possible on legitimate tasks.

% ============================================================
\section{RMU and the Forget--Retain Principle}
% ============================================================

\subsection*{Mathematical explanation}

RMU contributes an important mathematical pattern to NPS.

Instead of optimizing only for removal, it simultaneously:

\begin{enumerate}
    \item disrupts representations on hazardous data, and
    \item preserves representations on benign data.
\end{enumerate}

The general principle is:

\begin{equation}
\text{remove the dangerous computation}
+
\text{retain useful computation}.
\end{equation}

This is directly relevant to a firewall because a safe transformation that
destroys the model's capabilities is not sufficient.

However, RMU is not equivalent to NPS. RMU modifies the model to remove
hazardous knowledge, whereas NPS seeks an inference-time protected boundary.
Moreover, the WMDP paper reports that finetuning can recover the unlearned
capability in its open-source threat model. Thus representation disruption
does not automatically imply permanent privilege separation.

\subsection*{Theoretical definition}

A generic forget--retain objective is

\begin{equation}
\mathcal L
=
\mathcal L_{\mathrm{forget}}
+
\alpha
\mathcal L_{\mathrm{retain}}.
\end{equation}

For RMU specifically, the forget objective pushes hazardous activations
toward a fixed random direction while the retain objective keeps benign
activations near the frozen model's activations.

NPS generalizes the principle into a constraint problem:

\begin{equation}
\min
\quad
\mathcal L_{\mathrm{cap}}
+
\lambda\mathcal L_{\mathrm{iso}}
\end{equation}

subject to a policy-safety condition.

The conceptual shift is

\begin{equation}
\boxed{
\text{representation removal}
\quad\rightarrow\quad
\text{representation protection}.
}
\end{equation}

% ============================================================
\section{Policy Invariance}
% ============================================================

\subsection*{Mathematical explanation}

A true firewall should not merely correct a state once. If later transformer
layers can immediately recreate the forbidden state, the correction is not a
security boundary.

This motivates the concept of an invariant set from dynamical-systems
theory.

If the policy state begins inside the safe region, every subsequent protected
transition should remain inside it.

The key idea is:

\begin{equation}
\text{safe now}
\Rightarrow
\text{safe after the next transition}.
\end{equation}

Repeated application then gives safety throughout the entire computation.

\subsection*{Theoretical definition}

Let

\begin{equation}
P_\ell\in\mathcal P
\end{equation}

be the policy state and let

\begin{equation}
\Psafe\subseteq\mathcal P.
\end{equation}

The safe set is invariant if

\begin{equation}
\boxed{
P_\ell\in\Psafe
\Longrightarrow
P_{\ell+1}\in\Psafe.
}
\end{equation}

\begin{proposition}[Policy-state invariance]
If
\begin{equation}
P_0\in\Psafe
\end{equation}
and
\begin{equation}
P_\ell\in\Psafe
\Longrightarrow
P_{\ell+1}\in\Psafe
\end{equation}
for every layer, then
\begin{equation}
\boxed{
P_\ell\in\Psafe,\qquad\forall\ell.
}
\end{equation}
\end{proposition}

\begin{proof}
By induction. The base case is $P_0\in\Psafe$. If $P_\ell\in\Psafe$, the
invariance condition implies $P_{\ell+1}\in\Psafe$. Therefore the property
holds for all layers.
\end{proof}

This is the core candidate security invariant of NPS.

% ============================================================
\section{Policy Violation as a Behavioral Quantity}
% ============================================================

\subsection*{Mathematical explanation}

An internal state is only useful if it connects to the behavior we actually
care about.

A policy violation should therefore be represented as a measurable function
of the model's output or action. The internal firewall then has a chain of
meaning:

\begin{equation}
P
\rightarrow
\text{allowed/forbidden computation}
\rightarrow
Y
\rightarrow
\text{policy violation}.
\end{equation}

This prevents us from protecting an arbitrary latent feature that has no
reliable relationship to safety.

\subsection*{Theoretical definition}

Let

\begin{equation}
V(y)\in[0,1]
\end{equation}

be a policy-violation function.

For input $x$, define expected violation

\begin{equation}
R(x)
=
\E_{y\sim p_\theta(\cdot\mid x)}
[V(y)].
\end{equation}

For an adversarial input domain $\X_{\mathrm{adv}}$, the desired worst-case
security bound is

\begin{equation}
\boxed{
\sup_{x\in\X_{\mathrm{adv}}}
R_F(x)
\leq\epsilon.
}
\end{equation}

% ============================================================
\section{Adversarial Robustness and Recovery}
% ============================================================

\subsection*{Mathematical explanation}

A firewall is not secure merely because ordinary prompts fail to bypass it.
An attacker can optimize against the mechanism.

This creates a second-order question:

\begin{quote}
Can the attacker discover a different route through the model that avoids the
protected representation?
\end{quote}

RMU's evaluations illustrate the importance of this distinction: the authors
tested probes and adversarial attacks, but also found that open-source
finetuning could recover the unlearned capability. The lesson for NPS is that
``not currently observable'' and ``not recoverable under the threat model''
are different claims.

\subsection*{Theoretical definition}

Let $\mathcal A$ be an attacker class. Define

\begin{equation}
R_{\mathrm{adv}}(M_F)
=
\sup_{A\in\mathcal A}
V(A(M_F,x)).
\end{equation}

Require

\begin{equation}
\boxed{
R_{\mathrm{adv}}(M_F)\leq\rho.
}
\end{equation}

A stronger recovery formulation is

\begin{equation}
\sup_{A\in\mathcal A}
\Prob\left[
\text{recover forbidden behavior}
\right]
\leq\rho.
\end{equation}

The strength of this guarantee depends completely on the definition of
$\mathcal A$.

% ============================================================
\section{Representation Identifiability}
% ============================================================

\subsection*{Mathematical explanation}

Before protecting a policy state, we must establish that the candidate state
actually represents policy.

This is harder than obtaining high probe accuracy.

A classifier may exploit:

\begin{itemize}
    \item vocabulary,
    \item topic,
    \item prompt template,
    \item refusal style,
    \item dataset artifacts,
    \item or other variables correlated with the label.
\end{itemize}

The representation is interesting only if it continues to carry policy
information under changes that break those superficial correlations.

The backdoor paper reinforces this general lesson: robust conclusions require
explicit assumptions about what distinguishes the target feature from
naturally occurring features.

\subsection*{Theoretical definition}

Let $\Pi$ denote the underlying policy variable, $P$ the candidate latent
representation, and $X_{\mathrm{surface}}$ surface-level prompt properties.

A minimal information condition is

\begin{equation}
\boxed{
I(P;\Pi)>0.
}
\end{equation}

A desired specificity condition is

\begin{equation}
\boxed{
I(P;X_{\mathrm{surface}}\mid\Pi)\approx0.
}
\end{equation}

These conditions are not sufficient for causality. They motivate
intervention experiments.

% ============================================================
\section{Policy--User Independence}
% ============================================================

\subsection*{Mathematical explanation}

Suppose we find a representation that predicts policy perfectly but also
contains a large amount of user-controlled information.

That representation may be dangerous to protect because the attacker may be
able to use that extra information as a control channel.

Therefore NPS needs a stronger property than policy predictability:

\begin{equation}
\text{policy information}
\quad\text{without unnecessary attacker-controlled degrees of freedom}.
\end{equation}

This is the information-theoretic analogue of a privilege boundary.

\subsection*{Theoretical definition}

Let $U$ be user state and $P$ protected policy state. A desired conditional
independence property is

\begin{equation}
\boxed{
I(P;U\mid\Pi)\approx0.
}
\end{equation}

This should be interpreted carefully. It does not mean that the model cannot
understand the user's request. It means that, once the actual policy state is
fixed, the protected representation should contain little additional
user-controlled information.

% ============================================================
\section{The Backdoor Paper and the Need for Assumptions}
% ============================================================

\subsection*{Mathematical explanation}

The backdoor work provides a methodological warning that is directly relevant
to NPS.

There is no universal mathematical procedure for identifying ``the malicious
feature'' without assumptions. A malicious feature can be statistically
indistinguishable from a naturally occurring feature.

Their solution is to make an assumption tied to the effectiveness of the
attack---specifically, that the backdoor trigger is the strongest feature in
the relevant sense---and then derive a detection procedure under that
assumption.

NPS should adopt the same discipline.

We should not state:

\begin{equation}
\text{``there exists a universal policy representation''}
\end{equation}

without stating why such a representation should be identifiable.

Instead, we should formulate explicit structural assumptions and then derive
conditional guarantees.

\subsection*{Theoretical definition}

Let $\mathcal H$ denote the set of possible hypotheses about representation
structure and $\mathcal A$ the attacker class.

A theorem of the form

\begin{equation}
\mathcal H\land\mathcal A
\Longrightarrow
\text{security guarantee}
\end{equation}

is meaningful only when $\mathcal H$ and $\mathcal A$ are stated explicitly.

The backdoor paper provides an example of this methodology through assumptions
about feature strength and datamodel accuracy. \cite{backdoor}

For NPS, candidate assumptions may concern:

\begin{enumerate}
    \item existence of a stable policy representation,
    \item separability of policy from user-controlled state,
    \item bounded policy-state sensitivity,
    \item bounded attacker access,
    \item stability of protected transitions.
\end{enumerate}

These are hypotheses to test, not facts to assume silently.

% ============================================================
\section{Neural Firewall Security Definition}
% ============================================================

\subsection*{Mathematical explanation}

We can now combine the pieces.

A neural firewall should simultaneously provide:

\begin{enumerate}
    \item a meaningful policy representation,
    \item a protected safe region,
    \item limited user influence over that region,
    \item invariance of the region during computation,
    \item low policy violation under attack,
    \item and preservation of legitimate capabilities.
\end{enumerate}

No single scalar captures all of these properties. The firewall is therefore
best understood as a constrained dynamical system with a security invariant
and performance constraints.

\subsection*{Theoretical definition}

A candidate neural firewall consists of

\begin{equation}
\mathfrak F=
(M_F,\phi_P,\Psafe,\mathcal A,V,D),
\end{equation}

where:

\begin{itemize}
    \item $M_F$ is the protected model,
    \item $\phi_P$ extracts policy state,
    \item $\Psafe$ is the safe policy set,
    \item $\mathcal A$ is the attacker class,
    \item $V$ is the policy-violation function,
    \item $D$ measures capability deviation.
\end{itemize}

A strong NPS candidate should satisfy

\begin{align}
\text{Identifiability:}\qquad
&I(P;\Pi)>0,
\\
\text{User isolation:}\qquad
&I(P;U\mid\Pi)\approx0,
\\
\text{Local isolation:}\qquad
&
\left\|
\frac{\partial P'}{\partial U}
\right\|
\leq\epsilon,
\\
\text{Global isolation:}\qquad
&
\sup_{x'\in\A(x)}
d_P(P(x'),P(x))
\leq\eta,
\\
\text{Invariance:}\qquad
&
P_\ell\in\Psafe
\Rightarrow
P_{\ell+1}\in\Psafe,
\\
\text{Security:}\qquad
&
\sup_{x\in\X_{\mathrm{adv}}}R_F(x)
\leq\epsilon,
\\
\text{Capability:}\qquad
&
\E_{x\sim\X_{\mathrm{benign}}}D(p_F,p_M)
\leq\delta,
\\
\text{Recovery resistance:}\qquad
&
\sup_{A\in\mathcal A}R_{\mathrm{adv}}(M_F)
\leq\rho.
\end{align}

This is a target specification, not an established theorem.

% ============================================================
\section{The Neural Firewall Optimization Problem}
% ============================================================

\subsection*{Mathematical explanation}

The previous section gives constraints. We also need a constructive problem:
how should a firewall be learned or engineered?

The natural formulation is constrained optimization.

We want to minimize the amount of capability lost and the amount of policy
state exposed to user control, subject to a hard or soft security constraint.

This formulation makes the trade-offs explicit instead of hiding them in one
``safety score.''

\subsection*{Theoretical definition}

Let $\theta_F$ denote firewall parameters and $\phi$ the policy representation
map. Define

\begin{equation}
\boxed{
\begin{aligned}
\min_{\theta_F,\phi}\quad&
\mathcal L_{\mathrm{cap}}
+
\lambda\mathcal L_{\mathrm{iso}}
\\[3pt]
\text{s.t.}\quad&
\sup_{x\in\X_{\mathrm{adv}}}R_F(x)\leq\epsilon,
\\
&
\sup_{x'\in\A(x)}
d_P(P_F(x'),P_F(x))
\leq\eta,
\\
&
\sup_{A\in\mathcal A}
R_{\mathrm{adv}}(M_F)
\leq\rho.
\end{aligned}
}
\end{equation}

The objective may be expanded as

\begin{equation}
\mathcal L_{\mathrm{cap}}
=
\E_{x\sim\X_{\mathrm{benign}}}
D\!\left(
p_F(\cdot\mid x),
p_M(\cdot\mid x)
\right),
\end{equation}

with $\mathcal L_{\mathrm{iso}}$ representing a chosen measure of
user-to-policy leakage or sensitivity.

This is a research formulation rather than a solved optimization problem.

% ============================================================
\section{What Has Been Established Versus What NPS Must Prove}
% ============================================================

\subsection*{Mathematical explanation}

The four papers provide pieces of the theory, but they do not prove the NPS
claim.

RepE supports the idea that high-level concepts can be read and manipulated
in representation space and explicitly recommends converging evidence beyond
correlation.

ITI demonstrates inference-time causal intervention on a sparse set of
attention heads, while also exposing the truthfulness--helpfulness trade-off.

WMDP/RMU demonstrates representation-based suppression of hazardous knowledge
with retain objectives and evaluates recovery through probes and attacks, but
also shows that finetuning can recover the capability in an open-source
setting.

The backdoor paper demonstrates that robust feature identification requires
explicit assumptions.

Together, these sources motivate the framework but leave the central NPS
problem open:

\begin{equation}
\boxed{
\text{Can policy computation itself be protected from user-controlled
computation?}
}
\end{equation}

\subsection*{Theoretical definition}

The current state of knowledge can be summarized as:

\begin{center}
\begin{tabular}{ll}
\toprule
Question & Status \\
\midrule
Meaningful latent concepts exist & Supported \\
Concepts can be probed & Supported \\
Some concepts can be causally manipulated & Supported \\
Representation-level suppression is possible & Supported \\
Robustness requires explicit threat assumptions & Supported \\
Universal policy representation exists & Open \\
User/capability/policy separation exists & Open \\
Protected policy state exists & Open \\
Policy safe-set invariance is achievable & Open \\
Capability-preserving privilege separation is achievable & Open \\
\bottomrule
\end{tabular}
\end{center}

% ============================================================
\section{Research Program Derived from the Theory}
% ============================================================

\subsection*{Mathematical explanation}

The mathematics gives us a sequence of increasingly strong claims. This
prevents NPS from jumping directly from a successful probe to a claim of
security.

The progression is:

\begin{equation}
\boxed{
\text{representation}
\rightarrow
\text{causality}
\rightarrow
\text{isolation}
\rightarrow
\text{invariance}
\rightarrow
\text{robust security}.
}
\end{equation}

Each step adds a stronger requirement.

\subsection*{Theoretical definition}

\paragraph{Stage 1: Representation}

Establish

\begin{equation}
h\rightarrow P.
\end{equation}

\paragraph{Stage 2: Causality}

Establish

\begin{equation}
do(P\leftarrow P')
\Rightarrow
Y\text{ changes}.
\end{equation}

\paragraph{Stage 3: Isolation}

Establish

\begin{equation}
\sup_{x'\in\A(x)}
d_P(P(x'),P(x))
\leq\eta.
\end{equation}

\paragraph{Stage 4: Invariance}

Establish

\begin{equation}
P_\ell\in\Psafe
\Rightarrow
P_{\ell+1}\in\Psafe.
\end{equation}

\paragraph{Stage 5: Robustness}

Establish

\begin{equation}
\sup_{A\in\mathcal A}
V(A(M_F))
\leq\epsilon.
\end{equation}

% ============================================================
\section{The Central Theorem We Ultimately Want}
% ============================================================

\subsection*{Mathematical explanation}

The ultimate goal is not a theorem saying that a probe is accurate. It is a
theorem saying that a protected policy state is an invariant under an
explicitly defined attacker model.

Conceptually:

\begin{equation}
\text{safe initial policy state}
+
\text{protected transition}
\Rightarrow
\text{safe policy state for all computation}.
\end{equation}

If policy safety can also be shown to imply bounded behavioral violation,
then the latent invariant becomes a behavioral security guarantee.

This would be the mathematical core of a neural firewall.

\subsection*{Theoretical definition}

A target theorem has the form:

\begin{assumption}
For every allowed attacker input $x\in\X_{\mathrm{adv}}$, the protected
transition satisfies

\begin{equation}
P_\ell\in\Psafe
\Longrightarrow
P_{\ell+1}\in\Psafe.
\end{equation}

Furthermore, every state in $\Psafe$ induces bounded policy violation:

\begin{equation}
P_\ell\in\Psafe
\Longrightarrow
R_F(x)\leq\epsilon.
\end{equation}
\end{assumption}

\begin{proposition}[Target NPS security invariant]
Under the above assumptions, if

\begin{equation}
P_0\in\Psafe,
\end{equation}

then

\begin{equation}
P_\ell\in\Psafe
\qquad
\forall\ell,
\end{equation}

and consequently

\begin{equation}
R_F(x)\leq\epsilon
\qquad
\forall x\in\X_{\mathrm{adv}}.
\end{equation}
\end{proposition}

\begin{proof}
The invariant follows by induction from the protected transition condition.
The behavioral bound follows from the assumption that every state in
$\Psafe$ induces violation at most $\epsilon$.
\end{proof}

The difficult research problem is therefore not the induction itself. It is
constructing and validating the representation, safe set, transition
restriction, and attacker assumptions required for the theorem to apply.

% ============================================================
\section{Current Theoretical Boundary}
% ============================================================

\subsection*{Mathematical explanation}

At this point we have a coherent mathematical language for NPS, but not yet
a proof that real transformers satisfy it.

The following statements remain hypotheses:

\begin{itemize}
    \item a universal policy representation exists;
    \item policy can be separated from user-controlled computation;
    \item policy state has a stable geometry across prompts and tasks;
    \item the safe set can be characterized reliably;
    \item a protected transition can be implemented without unacceptable
          capability loss;
    \item the resulting invariant survives a sufficiently strong attacker.
\end{itemize}

This boundary is important. The mathematical framework should guide the
experiments rather than be mistaken for evidence that the desired structure
already exists.

\subsection*{Theoretical definition}

NPS is currently a conditional theory:

\begin{equation}
\boxed{
\text{If a stable, isolatable policy state exists,}
}
\end{equation}

then we seek to construct

\begin{equation}
\boxed{
\text{a protected invariant over that state}
}
\end{equation}

such that

\begin{equation}
\boxed{
\text{attacker-controlled input cannot induce unauthorized policy
violation within the stated threat model.}
}
\end{equation}

The next experiments therefore determine whether the premises of this theory
are empirically satisfied.

% ============================================================
\section{Conclusion}
% ============================================================

\subsection*{Mathematical explanation}

The current theory can be summarized as a progression from representation
science to security theory.

Representation Engineering motivates the search for meaningful latent
variables. ITI shows that some latent variables can be manipulated at
inference time. RMU shows that representation-level modification can suppress
hazardous capability while retaining some useful behavior, but also exposes
recovery and capability-loss problems. The backdoor work establishes the need
for explicit assumptions when making robust feature-level claims.

NPS adds a stronger target:

\begin{equation}
\text{not merely ``find and steer a safety direction,''}
\end{equation}

but

\begin{equation}
\boxed{
\text{identify}
\rightarrow
\text{isolate}
\rightarrow
\text{protect}
\rightarrow
\text{prove invariance}.
}
\end{equation}

The proposed neural firewall is therefore a constrained dynamical system
whose protected policy state acts as privileged internal state.

\subsection*{Theoretical definition}

The current mathematical abstraction is

\begin{equation}
\boxed{
\text{Neural Firewall}
=
\text{Protected Policy State}
+
\text{Constrained Transition}
+
\text{Capability Preservation}
+
\text{Adversarial Robustness}.
}
\end{equation}

The central open problem is

\begin{equation}
\boxed{
\exists\,\phi,\Psafe,\mathcal F
\quad\text{such that}\quad
\begin{cases}
P=\phi_P(h)\text{ is policy-relevant},\\
U\not\rightarrow P\text{ under }\mathcal A,\\
P_\ell\in\Psafe\Rightarrow P_{\ell+1}\in\Psafe,\\
R_F(x)\leq\epsilon,\\
\mathcal L_{\mathrm{cap}}\leq\delta.
\end{cases}
}
\end{equation}

That existence question is the mathematical heart of NPS.

% ============================================================
\begin{thebibliography}{9}

\bibitem{repe}
Andy Zou et al.
\newblock \emph{Representation Engineering: A Top-Down Approach to AI
Transparency}.
\newblock arXiv:2310.01405, v4, 2025.

\bibitem{iti}
Kenneth Li et al.
\newblock \emph{Inference-Time Intervention: Eliciting Truthful Answers from
a Language Model}.
\newblock arXiv:2306.03341, v6.

\bibitem{wmdp}
Nathaniel Li et al.
\newblock \emph{The WMDP Benchmark: Measuring and Reducing Malicious Use With
Unlearning}.
\newblock arXiv:2403.03218, v7, 2024.

\bibitem{backdoor}
Khaddaj et al.
\newblock \emph{Rethinking Backdoor Attacks}.
\newblock arXiv:2307.10163, v1.

\bibitem{nps}
NPS Research Charter.
\newblock \emph{Neural Privilege Separation: Research Vision and Roadmap}.
\newblock Internal project document.

\end{thebibliography}

\end{document}
